\documentclass[11pt]{article}
\usepackage[utf8]{inputenc}
\usepackage[T1]{fontenc}
\usepackage[french]{babel}
\usepackage[margin=2.5cm]{geometry}
\usepackage{hyperref}

\title{Journal de travail --- Dashboard TC1.1--1.5 \& flag \texttt{real\_flag}}
\author{DeepEdgeBenchmark}
\date{20 juillet 2026}

\begin{document}
\maketitle

\section{Contexte}
Point de départ : le dashboard signalait des données manquantes pour ARIMA
(``seulement 2 jours sur 9 avec signal''). Vérification faite directement en
base (\texttt{validation/tracking.db}) plutôt que de faire confiance au
diagnostic transmis.

\section{Diagnostic du ``trou de données''}
\begin{itemize}
  \item La base locale était périmée (dernier \texttt{cutoff\_date} au 15/07,
        alors que des runs plus récents existaient déjà sur \texttt{origin/main}).
  \item Après \texttt{git pull}, les prédictions du 16 au 19/07 sont apparues,
        mais leurs \texttt{sim\_trades} n'étaient pas encore générés : le cron
        \texttt{evaluate\_daily} n'avait simplement pas encore tourné après ce
        backfill (fait par Maeva Cavaleri le jour même).
  \item Les chiffres ``2 jours sur 9 / 6 jours sur 9'' avancés dans l'échange
        précédent ne correspondaient à aucune donnée réelle en base : ARIMA
        avait en fait des signaux sur 9 jours sur 9 (\texttt{bull\_calm\_d1} et
        \texttt{sideways\_d1}), simplement avec un taux de réussite directionnel
        plus faible que SARIMA -- une question de performance, pas de collecte.
  \item Conclusion : pas de bug à corriger, juste attendre le prochain passage
        du cron \texttt{evaluate-daily.yml} (GitHub Actions, 01h30 Paris).
\end{itemize}

\section{Extension du tableau Test Cases à tous les actifs}
Le tableau TC1.1--1.5 du dashboard (\texttt{model\_artifacts/generate\_dashboard.py})
n'affichait que BTC-USD (\texttt{SIM\_TRADES\_ASSETS = ["BTC-USD"]}), alors que
le rendu JS était déjà générique par actif.
\begin{itemize}
  \item \texttt{SIM\_TRADES\_ASSETS} tire désormais la liste depuis
        \texttt{calibration.regime.assets.ASSETS} (source unique de vérité),
        avec repli sur la liste connue si l'import échoue.
  \item Vérifié : les 5 actifs (BTC-USD, ETH-USD, SPY, ZN=F, TLT) ont bien
        leurs lignes dans le dashboard régénéré.
  \item Commit \texttt{98aaff5}, poussé sur \texttt{main} --- déclenche le
        redéploiement automatique de \texttt{deploy-pages.yml}.
\end{itemize}

\section{Nouveau flag \texttt{real\_flag} (vraie vs fausse prédiction)}

\subsection{Le piège identifié avant de coder}
La colonne \texttt{source} (\texttt{'live'}/\texttt{'oos'}) porte déjà un sens
technique (provenance : pipeline de prod vs backtest), utilisé comme clé de
jointure et d'unicité dans une bonne dizaine d'endroits (\texttt{sim\_trades},
index unique OOS, \texttt{evaluate\_pending}, dédup, gate stress de
\texttt{sideways\_gated\_d1}\ldots). Un historique git (\texttt{f60d71f} puis
\texttt{d675188}, revert le jour même) montrait qu'une tentative précédente de
réécrire \texttt{source} selon vraie/fausse prédiction avait cassé la jointure
\texttt{daily\_detail()} et rendu les signaux TC invisibles --- exactement le
symptôme diagnostiqué en section 2.

\subsection{Conception retenue}
\begin{itemize}
  \item Nouvelle colonne \textbf{additive} \texttt{predictions.real\_flag}
        (\texttt{'live'} = vraie prédiction, \texttt{'oos'} = fausse), jamais
        utilisée comme clé de jointure ou d'unicité.
  \item Règle unique de calcul : \texttt{tracking\_db.compute\_real\_flag(model,
        cutoff\_date)} --- reprend exactement le seuil déjà utilisé côté JS
        (\texttt{isRealPrediction}), qui devient la seule source de vérité.
  \item Migration auto-réparatrice dans \texttt{init\_db()} : ajoute la colonne
        si absente et backfille toute ligne existante à chaque appel.
  \item Calcul également fait à l'insertion (\texttt{save\_prediction},
        \texttt{insert\_oos\_predictions}) en défense en profondeur.
  \item Propagation jusqu'au dashboard : vue \texttt{all\_predictions}
        $\rightarrow$ \texttt{daily\_detail()} $\rightarrow$
        \texttt{collect\_sim\_trades\_daily()} $\rightarrow$ JS
        (\texttt{isRealPrediction(row) => row.real\_flag === 'live'}).
\end{itemize}

\subsection{Vérifications avant commit}
\begin{itemize}
  \item Tests ciblés : \texttt{pytest validation/test\_tracking\_db.py
        validation/test\_sim\_trades.py -q} $\rightarrow$ 179 tests passés
        (dont 5 nouveaux tests couvrant la migration et le cas backfill OOS).
  \item Migration testée sur une copie de la base : aucune perte de ligne,
        0 valeur \texttt{NULL} restante.
  \item Spot-check exact sur les dates de backfill (08, 11, 13/07) :
        \texttt{source='oos'} et \texttt{real\_flag='live'} coexistent bien sur
        la même ligne --- le cas que \texttt{source} seul ne pouvait pas
        représenter.
  \item Dashboard régénéré : mêmes comptes de lignes par actif qu'avant le
        changement (886/889/609/614/609) --- aucune régression du type
        \texttt{d675188}.
\end{itemize}

\subsection{Commit}
\texttt{71f3179} --- \emph{feat(db): flag real\_flag ('live'/'oos') pour
vraie/fausse prédiction}, mergé avec un commit concurrent de Maeva
(\texttt{documentation/oos\_vs\_live.pdf}, sans conflit), poussé sur
\texttt{main} sous \texttt{d2a110e}.

\section{Fichiers modifiés aujourd'hui}
\begin{itemize}
  \item \texttt{validation/tracking\_db.py} --- schéma, migration, \texttt{compute\_real\_flag}.
  \item \texttt{validation/sim\_trades.py} --- vue \texttt{all\_predictions}, \texttt{daily\_detail}, \texttt{insert\_oos\_predictions}.
  \item \texttt{model\_artifacts/generate\_dashboard.py} --- \texttt{SIM\_TRADES\_ASSETS}, \texttt{collect\_sim\_trades\_daily}, JS \texttt{isRealPrediction}.
  \item \texttt{validation/test\_tracking\_db.py}, \texttt{validation/test\_sim\_trades.py} --- nouveaux tests.
  \item \texttt{validation/tracking.db} --- backfill réel de \texttt{real\_flag}.
\end{itemize}

\section{Suite}
Le cron \texttt{evaluate-daily.yml} (01h30 Paris) doit résoudre les
prédictions du 16--19/07 et générer leurs \texttt{sim\_trades} manquants ---
rien d'autre à faire de ce côté. Le dashboard en ligne est à jour suite aux
deux pushs du jour.

\end{document}
